\documentclass[11pt]{article}

\usepackage[a4paper,margin=1in]{geometry}
\usepackage{amsmath,amssymb,braket}
\usepackage{physics}
\usepackage{hyperref}
\usepackage{enumitem}
\usepackage{array}

\title{From Thermodynamic Neurons to Floquet-Buffered Non-Markovian Logic Gates}
\author{Implementation and Theory Bridge}
\date{\today}

\begin{document}
\maketitle

\begin{abstract}
This note explains how the current codebase connects to the proposed research idea.  The
baseline thermodynamic-neuron paper is first reconstructed in the Markovian reset model.
The extension then asks what changes when the bath is structured, has memory, and is not
directly attached to the logical qubit.  The proposed answer is to insert a periodically driven
Floquet buffer between the logical subsystem and the environment.  The buffer acts as a
dynamical filter and isolation layer.  Since direct simulation of a many-qubit system plus
structured reservoirs grows exponentially, the environmental memory is represented as a
tensor network, specifically a process tensor or PT-MPO.  The operational output is not an
abstract current alone, but a measurable logical distinguishability of output states.
\end{abstract}

\section{What Has Already Been Reconstructed}

The first completed part of the project is the Markovian baseline thermodynamic neuron.
This corresponds to the weak-coupling, reset-model picture of the baseline paper.  In this
model, logical inputs are encoded as inverse temperatures of thermal reservoirs, and the
output is encoded in the final inverse temperature of a finite reservoir.

The basic thermal occupation of a qubit with energy gap $\epsilon$ is
\begin{equation}
    g(\beta \epsilon)
    =
    \frac{1}{1 + e^{\beta \epsilon}}.
\end{equation}

For the three-qubit collector of the NOT gate, the virtual inverse temperature is
\begin{equation}
    \beta_v
    =
    \frac{\epsilon_0}{\epsilon_z}\beta_0
    -
    \frac{\epsilon_1}{\epsilon_z}\beta_1,
    \qquad
    \epsilon_z = \epsilon_0 - \epsilon_1.
\end{equation}

The reset-model collector and modulator currents are
\begin{align}
    j_C
    &=
    \mu \epsilon_z
    \left[
        g_z(\beta_z) - g_z(\beta_v)
    \right],
    \\
    j_M
    &=
    \mu' \epsilon_z
    \left[
        g_z(\beta_z) - g_z(\beta_r)
    \right].
\end{align}
Here $g_z(\beta) = g(\beta \epsilon_z)$.  The finite output reservoir evolves as
\begin{equation}
    \dot{\beta}_z(t)
    =
    -
    \frac{\beta_z(t)^2}{C}
    \left[
        j_C + j_M
    \right].
\end{equation}

The code in \texttt{Scripts/baseline\_paper\_reconstruction/} implements this model and
generates the baseline NOT, NOR, 3-majority, and XOR responses.  This is the reference point:
every extension must be compared against it.

\section{Why The Baseline Is Not The Final Theory}

The baseline model is clean because the bath is memoryless.  The system state $\rho_S(t)$
evolves using only its current value:
\begin{equation}
    \dot{\rho}_S(t)
    =
    \mathcal{L}\rho_S(t).
\end{equation}
This is a Markovian GKSL-type picture.  It assumes weak system-bath coupling, rapid bath
correlation decay, and no long-lived information returning from the reservoir.

For a structured bath, this is no longer safe.  The bath may remember what the system did
earlier.  Then the reduced dynamics has a memory kernel:
\begin{equation}
    \dot{\rho}_S(t)
    =
    -i[H_S,\rho_S(t)]
    +
    \int_0^t
    \mathcal{K}(t-\tau)\rho_S(\tau)\,d\tau.
\end{equation}
This is the Nakajima--Zwanzig form.  It says, in plain language, that the present motion of the
system depends on its past.  The object $\mathcal{K}(t-\tau)$ is the memory kernel.  When
the bath memory is very short, the kernel becomes sharply peaked and one recovers the
ordinary Markovian equation.  When the memory is long, the Markovian approximation may
miss revivals, backflow, and modified relaxation.

This is why the project cannot jump directly from the reset-model baseline to a strong claim
about non-Markovian thermodynamic gates.  First one must show that the non-Markovian
solver is stable and that the chosen observable is meaningful.

\section{The Floquet Buffer Idea}

The proposed architecture inserts a driven intermediate system between the logical subsystem
and the bath:
\begin{equation}
    S \longleftrightarrow F(t) \longleftrightarrow B.
\end{equation}
The direct architecture is
\begin{equation}
    S \longleftrightarrow B.
\end{equation}

Here $S$ is the logical subsystem, $F(t)$ is the periodically driven Floquet buffer, and $B$ is
the bath.  A minimal buffer Hamiltonian is
\begin{equation}
    H_F(t)
    =
    \frac{\epsilon_F}{2}\sigma_z
    +
    A\cos(\Omega t)\sigma_x,
\end{equation}
with period
\begin{equation}
    T = \frac{2\pi}{\Omega}.
\end{equation}
The logical qubit-buffer interaction can be chosen as
\begin{equation}
    H_{SF}
    =
    g_{SF}\,
    \sigma_x^{(S)}\otimes\sigma_x^{(F)}.
\end{equation}

The complete minimal bridge Hamiltonian is therefore
\begin{equation}
    H(t)
    =
    H_S
    +
    H_F(t)
    +
    H_{SF}
    +
    H_B
    +
    H_{FB}.
\end{equation}

The physical role of the buffer is simple:
\begin{itemize}[leftmargin=2em]
    \item it prevents the logical qubit from touching the bath directly;
    \item it creates drive-controlled sidebands at frequencies shifted by multiples of $\Omega$;
    \item it can suppress unwanted transitions and enhance desired transitions;
    \item it gives a possible gain/isolation stage, analogous in spirit to a CMOS inverter.
\end{itemize}

In a weak Markovian treatment of the buffer-bath contact, one obtains a time-periodic
effective generator
\begin{equation}
    \dot{\rho}_S(t)
    =
    \mathcal{L}_{\mathrm{eff}}(t)\rho_S(t),
    \qquad
    \mathcal{L}_{\mathrm{eff}}(t+T)
    =
    \mathcal{L}_{\mathrm{eff}}(t).
\end{equation}
The transition rates depend on bath correlations evaluated at shifted frequencies
\begin{equation}
    \omega + n\Omega,
\end{equation}
where $n\in\mathbb{Z}$.  This is the origin of Floquet filtering.

\section{What The Current Floquet Code Does}

The code in \texttt{Scripts/floquet\_buffer\_extension/} implements the first bridge model:
\begin{equation}
    S \longleftrightarrow F(t) \longleftrightarrow B.
\end{equation}
It compares this against the direct baseline
\begin{equation}
    S \longleftrightarrow B.
\end{equation}

For two different initial logical states, $\rho_S^{(0)}(0)$ and $\rho_S^{(1)}(0)$, the code evolves
the output states and computes the trace distance
\begin{equation}
    D(t)
    =
    \frac{1}{2}
    \norm{
        \rho_S^{(0)}(t)
        -
        \rho_S^{(1)}(t)
    }_1.
\end{equation}
This is the operational distinguishability of the two logical outputs.  If $D=1$, the states are
perfectly distinguishable.  If $D=0$, they are indistinguishable.

The code also estimates the work injected by the drive through
\begin{equation}
    W_{\mathrm{drive}}
    =
    \int_0^{t_f}
    \Tr\left[
        \rho(t)
        \frac{\partial H_F(t)}{\partial t}
    \right]dt.
\end{equation}
For
\begin{equation}
    H_F(t)
    =
    \frac{\epsilon_F}{2}\sigma_z
    +
    A\cos(\Omega t)\sigma_x,
\end{equation}
we have
\begin{equation}
    \frac{\partial H_F(t)}{\partial t}
    =
    -
    A\Omega\sin(\Omega t)\sigma_x.
\end{equation}

The first default run gives
\begin{align}
    D_{\mathrm{direct}}(t_f)
    &\approx
    0.135,
    \\
    D_{\mathrm{buffer}}(t_f)
    &\approx
    0.182,
    \\
    W_{\mathrm{drive}}
    &\approx
    0.460.
\end{align}
Thus the first parameter set shows the intended trade-off:
\begin{equation}
    \text{better output distinguishability}
    \quad
    \text{at the cost of positive drive work.}
\end{equation}

\section{Why Tensor Networks Are Needed}

If one explicitly simulates the system plus bath, the Hilbert space grows exponentially.  Suppose
the system has dimension $d_S$ and the bath is represented by $N$ local sites of dimension $d$.
The full state dimension is
\begin{equation}
    d_S d^N.
\end{equation}
This is impossible for large $N$.

For non-Markovian open systems, the difficulty can also be seen in time.  A bath with memory
couples the system at time $t_k$ to its earlier states at times $t_j$.  Naively, storing all multi-time
correlations scales exponentially with the number of time steps.

The process tensor solves this by representing the environment's influence as a tensor network:
\begin{equation}
    \Upsilon_{n:0}
    =
    \sum_{\{\alpha_k,\beta_k\}}
    \Tr\left[
        B^{[n-1]}_{\alpha_{n-1},\beta_{n-1}}
        \cdots
        B^{[0]}_{\alpha_0,\beta_0}
    \right]
    \ket{\alpha_{n-1}\cdots\alpha_0}
    \bra{\beta_{n-1}\cdots\beta_0}.
\end{equation}
This is a matrix product operator along the time direction.  Its internal bond dimension
$\chi$ measures how much memory must be carried from the past to the future.

If the bath is Markovian, then
\begin{equation}
    \chi = 1.
\end{equation}
If the bath is non-Markovian but compressible, then
\begin{equation}
    1 < \chi \ll d^{2n}.
\end{equation}
This is the key reason tensor networks prevent exponential blow-up: the full memory object
is not stored exactly as an exponentially large tensor; it is compressed by singular value
decomposition while keeping the physically important temporal correlations.

\section{Single Example: From Nakajima--Zwanzig To Output Logic}

Consider one logical qubit coupled to a structured bath through a Floquet buffer:
\begin{equation}
    S \longleftrightarrow F(t) \longleftrightarrow B.
\end{equation}

\subsection*{Step 1: Exact Reduced Dynamics}

The full density matrix obeys
\begin{equation}
    \dot{\rho}_{SFB}(t)
    =
    -i[H_{SFB}(t),\rho_{SFB}(t)].
\end{equation}
The reduced logical state is
\begin{equation}
    \rho_S(t)
    =
    \Tr_{FB}\rho_{SFB}(t).
\end{equation}

\subsection*{Step 2: Memory Kernel View}

Tracing out the buffer and bath gives a memory equation
\begin{equation}
    \dot{\rho}_S(t)
    =
    -i[H_S,\rho_S(t)]
    +
    \int_0^t
    \mathcal{K}_F(t,\tau)\rho_S(\tau)d\tau.
\end{equation}
The kernel now depends on the driven buffer.  By tuning $A$ and $\Omega$, one changes
$\mathcal{K}_F$ and therefore changes how information leaves or returns to the logical qubit.

\subsection*{Step 3: Tensor Network Representation}

Instead of explicitly constructing $\mathcal{K}_F$, the process-tensor method builds the
multi-time influence object $\Upsilon_{n:0}$.  The simulation then contracts the logical system
operations with this environment tensor network:
\begin{equation}
    \rho_S(t_n)
    =
    \Upsilon_{n:0}
    \star
    \left(
        \mathcal{A}_{n-1}
        \otimes
        \cdots
        \otimes
        \mathcal{A}_0
    \right),
\end{equation}
where $\star$ means tensor contraction and $\mathcal{A}_k$ are the system operations at each
time step.

\subsection*{Step 4: Logical Output}

For a binary gate, prepare two input states and evolve them:
\begin{equation}
    \rho_S^{(0)}(0)
    \rightarrow
    \rho_S^{(0)}(t_f),
    \qquad
    \rho_S^{(1)}(0)
    \rightarrow
    \rho_S^{(1)}(t_f).
\end{equation}
The logical distinguishability is
\begin{equation}
    D_{\mathrm{out}}
    =
    \frac{1}{2}
    \norm{
        \rho_S^{(0)}(t_f)
        -
        \rho_S^{(1)}(t_f)
    }_1.
\end{equation}
The Floquet buffer is useful only if it improves $D_{\mathrm{out}}$ or improves robustness at an
acceptable drive cost.

\section{How The Codebase Now Maps To The Theory}

\begin{center}
\begin{tabular}{>{\raggedright\arraybackslash}p{0.24\linewidth}
                >{\raggedright\arraybackslash}p{0.36\linewidth}
                >{\raggedright\arraybackslash}p{0.28\linewidth}}
\hline
Theory object & Code location & Purpose \\
\hline
Baseline reset model
& baseline reconstruction folder
& Markovian thermodynamic neuron \\
Structured bath benchmark
& non-Markovian extension folder
& OQuPy/PT-MPO calibration \\
Floquet buffer bridge
& Floquet-buffer extension folder
& Test driven isolation layer \\
Convergence scans
& non-Markovian outputs folder
& Verify numerical stability \\
\hline
\end{tabular}
\end{center}

\section{Next Scientific Step}

The next step is not to claim that the full thermodynamic gate has been solved.  The next step is
to connect the Floquet-buffer bridge to the non-Markovian machinery:
\begin{equation}
    S \longleftrightarrow F(t) \longleftrightarrow B_{\mathrm{structured}},
\end{equation}
with the structured bath represented by a process tensor or chain mapping.  The first target
should remain a single logical qubit.  Only after this is stable should the full three-qubit NOT
or NOR thermodynamic neuron be attempted.

\end{document}
